% docs/slides/preamble.tex -- 五份 deck 共享
\documentclass[aspectratio=169,11pt]{beamer}

% ==== 中文与字体 ====
\usepackage[UTF8,fontset=none]{ctex}
\usepackage{fontspec}

\IfFontExistsTF{PingFang SC}{%
  \setCJKmainfont{PingFang SC}[BoldFont=PingFang SC,ItalicFont=PingFang SC,AutoFakeSlant=0.15]
  \setCJKsansfont{PingFang SC}[ItalicFont=PingFang SC,AutoFakeSlant=0.15]
  \setCJKmonofont{PingFang SC}
}{%
  \IfFontExistsTF{Source Han Sans SC}{%
    \setCJKmainfont{Source Han Sans SC}
    \setCJKsansfont{Source Han Sans SC}
  }{%
    \IfFontExistsTF{Noto Sans CJK SC}{%
      \setCJKmainfont{Noto Sans CJK SC}
      \setCJKsansfont{Noto Sans CJK SC}
    }{%
      \setCJKmainfont{Heiti SC}
      \setCJKsansfont{Heiti SC}
    }
  }
}

\IfFontExistsTF{Helvetica Neue}{\setmainfont{Helvetica Neue}}{%
  \IfFontExistsTF{Inter}{\setmainfont{Inter}}{}}
\IfFontExistsTF{JetBrains Mono}{\setmonofont{JetBrains Mono}[Scale=0.85]}%
  {\setmonofont{Menlo}[Scale=0.85]}

% ==== 颜色 ====
\usepackage{xcolor}
\definecolor{cMain}{HTML}{1F4E79}
\definecolor{cAccent}{HTML}{E07A1F}
\definecolor{cGray}{HTML}{595959}
\definecolor{cMute}{HTML}{F4F4F4}
\definecolor{cOK}{HTML}{2E7D32}
\definecolor{cBad}{HTML}{B23A3A}
\definecolor{cWarn}{HTML}{C8A300}

% ==== 主题 ====
\usepackage{tcolorbox}
\usetheme{metropolis}
\metroset{progressbar=frametitle,numbering=fraction,block=fill}
\setbeamercolor{frametitle}{bg=cMain,fg=white}
\setbeamercolor{progress bar}{fg=cAccent,bg=cMute}
\setbeamercolor{alerted text}{fg=cAccent}
\setbeamercolor{title}{fg=cMain}

% metropolis 默认尝试 Fira Sans，未安装时 fallback 到 Latin Modern Sans。
% 这里在主题加载后显式覆盖，强制使用 macOS 上一定存在的 Helvetica Neue，
% 跟中文 PingFang 风格一致；找不到则按顺序尝试 Inter / Arial。
\IfFontExistsTF{Helvetica Neue}{%
  \setsansfont{Helvetica Neue}%
}{%
  \IfFontExistsTF{Inter}{\setsansfont{Inter}}{%
    \IfFontExistsTF{Arial}{\setsansfont{Arial}}{}}%
}

% ==== TikZ ====
\usepackage{tikz}
\usetikzlibrary{positioning,arrows.meta,fit,shapes.geometric,calc,%
  decorations.pathmorphing,backgrounds}

\tikzset{
  node-box/.style={draw=cMain,thick,rounded corners=2pt,
    minimum height=8mm,minimum width=18mm,align=center,font=\small},
  node-state/.style={draw=cMain,thick,circle,minimum size=14mm,
    align=center,font=\small},
  node-ok/.style={draw=cOK,thick,fill=cOK!10},
  node-warn/.style={draw=cWarn,thick,fill=cWarn!15},
  node-bad/.style={draw=cBad,thick,fill=cBad!10},
  % 色盲友好的状态机三态：蓝 (正常/Closed) / 橙 (告警/Open) / 灰 (中间/HalfOpen)
  % 避免 red-green 二元区分，改用 hue + lightness 双通道。
  node-state-normal/.style={draw=cMain,thick,fill=cMain!12},
  node-state-alert/.style={draw=cAccent,thick,fill=cAccent!18},
  node-state-transition/.style={draw=cGray,thick,fill=cGray!18},
  arrow/.style={-{Stealth[length=2mm]},thick,draw=cMain},
  arrow-async/.style={-{Stealth[length=2mm]},thick,dashed,draw=cGray},
  arrow-bad/.style={-{Stealth[length=2mm]},thick,draw=cBad},
  arrow-alert/.style={-{Stealth[length=2mm]},thick,draw=cAccent},
  lifeline/.style={dashed,thick,draw=cGray},
}

% ==== 代码 ====
\usepackage{listings}
\lstdefinelanguage{Go}{
  morekeywords={break,case,chan,const,continue,default,defer,else,
    fallthrough,for,func,go,goto,if,import,interface,map,package,range,
    return,select,struct,switch,type,var,nil,true,false,iota},
  morekeywords=[2]{string,int,int64,uint,uint64,bool,byte,error,context},
  sensitive=true,
  morecomment=[l]//,morecomment=[s]{/*}{*/},
  morestring=[b]",morestring=[b]`,
}
\lstdefinelanguage{LuaRedis}{
  morekeywords={and,break,do,else,elseif,end,false,for,function,if,in,
    local,nil,not,or,repeat,return,then,true,until,while},
  morekeywords=[2]{redis,call,KEYS,ARGV,tonumber,HGET,HSET,GET,SET,
    DECR,DECRBY,INCR,INCRBY,EXPIRE,PEXPIRE,ZADD,ZCARD,ZREMRANGEBYSCORE,
    EVAL,EVALSHA},
  sensitive=true,
  morecomment=[l]{--},
  morestring=[b]",morestring=[b]',
}

\lstdefinestyle{base}{
  basicstyle=\ttfamily\scriptsize,
  numbers=left,numberstyle=\tiny\color{cGray},
  keywordstyle=\color{cMain}\bfseries,
  keywordstyle=[2]\color{cAccent},
  stringstyle=\color{cOK},
  commentstyle=\color{cGray}\itshape,
  backgroundcolor=\color{cMute},
  frame=single,framerule=0pt,
  showstringspaces=false,
  breaklines=true,
  tabsize=2,
  columns=fullflexible,
}
\lstset{style=base}

% ==== 三线表与附录页码 ====
\usepackage{booktabs}
\usepackage{appendixnumberbeamer}

% ==== 页脚来源标注 ====
\newcommand{\srcnote}[1]{{\color{cGray}\scriptsize\itshape #1}}

% 关键论点框
\newcommand{\keypoint}[1]{%
  \begin{tcolorbox}[colback=cMute,colframe=cMain,boxrule=0.5pt,arc=2pt]%
  #1\end{tcolorbox}}


\title{谁能做什么}
\subtitle{角色边界、双 token 续命、客服补券与黑产温床}
\author{gomall · 业务侧 deck}
\date{}

\begin{document}

\maketitle

\begin{frame}{今日大纲}
\tableofcontents
\end{frame}

% ============================================================
\section{第一性原理：鉴权解决的不是技术问题，是责任划分}
% ============================================================

\begin{frame}{一句话：鉴权 = 把"谁"和"能做什么"两件事分开}
\begin{itemize}
  \item "谁"（Authentication，认证）：你是不是张三？$\to$ 看 token。
  \item "能做什么"（Authorization，授权）：张三能不能改价？$\to$ 看角色。
  \item 这两件事必须分开，否则一旦合并：
  \begin{itemize}
    \item 客服认错人 = 真用户被代下单。
    \item 角色判错 = 普通用户能进 admin 后台。
  \end{itemize}
  \item gomall 的边界分两层：先 \texttt{AuthMiddleware}（是不是登录用户），再 \texttt{RequireRole}（是不是 admin）。
\end{itemize}
\srcnote{routes/router.go:45--46 authed group；:125--126 admin sub-group}
\end{frame}

% ============================================================
\section{角色边界：把 router 翻译成业务能力}
% ============================================================

\begin{frame}{gomall 现在只有两种角色}
\begin{tabular}{@{}lll@{}}
\toprule
角色 & 业务身份 & 路由前缀 \\
\midrule
匿名     & 没登录的访客         & \texttt{/api/v1/*}（白名单） \\
user     & 注册用户             & \texttt{/api/v1/*} + \texttt{authed.*} \\
admin    & 运营 / 客服 / 风控    & 以上 + \texttt{authed.admin.*} \\
\bottomrule
\end{tabular}
\vspace{3mm}
\small 没有"商家"，没有"运营组长"，没有"超级 admin"。\\
角色越少，越不容易出"谁可以批准谁"的扯皮。等业务真有分级再加，不要预先抽象。
\srcnote{repository/db/model/user.go:23 Role 字段；:30--31 RoleUser / RoleAdmin}
\end{frame}

\begin{frame}{三层墙：匿名 / user / admin}
\begin{tikzpicture}[node distance=6mm and 6mm,every node/.append style={align=center,font=\scriptsize,text width=30mm}]
  \node[node-box,node-state-normal] (open)  {匿名可见\\\texttt{product/show}\\\texttt{ping / register / login}};
  \node[node-box,node-state-transition,right=of open] (user) {登录后可用\\\texttt{orders/create}\\\texttt{coupon/claim, paydown}};
  \node[node-box,node-state-alert,right=of user] (admin) {仅 admin\\\texttt{users / promote}\\\texttt{search/backfill}};
  \draw[arrow]      (open) -- node[above,font=\tiny]{AuthMiddleware} (user);
  \draw[arrow-alert](user) -- node[above,font=\tiny]{RequireRole("admin")}(admin);
\end{tikzpicture}
\vspace{3mm}
\begin{itemize}
  \item \textbf{匿名层}：要能引流，所以商品详情 / 列表 / 类目都开放。
  \item \textbf{user 层}：涉及资金 / 个人数据，必须先确认"是不是你本人"。
  \item \textbf{admin 层}：动别人账户的能力 —— 必须二次校验角色。
\end{itemize}
\srcnote{routes/router.go:34--43 匿名；:45--46 authed；:125--131 admin}
\end{frame}

\begin{frame}{为什么这些接口在 authed 而那些不在}
\begin{tabular}{@{}lll@{}}
\toprule
接口 & 业务理由 & 路由位置 \\
\midrule
\texttt{product/show}    & 引流必须匿名可看  & router.go:39 \\
\texttt{user/login}      & 登录前怎么可能登录 & router.go:35 \\
\texttt{orders/create}   & 没人的订单是谁的？ & router.go:77（authed） \\
\texttt{coupon/claim}    & 没人怎么把券记到账户上？& router.go:70（authed） \\
\texttt{paydown}         & 资金类，最严      & router.go:96（authed） \\
\texttt{admin/users}     & 看全员数据，必须 admin & router.go:128 \\
\texttt{search/backfill} & 一次几千万条写 ES & router.go:130 \\
\bottomrule
\end{tabular}
\vspace{2mm}
\small "能不能匿名访问"不是开发拍的，是"业务上有没有公网引流诉求"决定的。
\end{frame}

% ============================================================
\section{双 token 续命：用户感受不到登录这件事}
% ============================================================

\begin{frame}{双 token 续期：用户感受不到登录这件事}
\begin{tikzpicture}[node distance=10mm and 12mm]
  \node[node-box] (app) {App 客户端};
  \node[node-box,right=of app,node-warn] (mw) {AuthMiddleware};
  \node[node-box,right=of mw] (svc) {业务 handler};
  \draw[arrow] (app) -- node[above,font=\tiny]{access + refresh} (mw);
  \draw[arrow] (mw) -- node[above,font=\tiny]{合法 $\to$ 透传} (svc);
  \draw[arrow-alert,bend left=35] (mw.north) to node[above,font=\tiny]{Set-Cookie 新 token 对} (app.north);
\end{tikzpicture}
\vspace{2mm}
\begin{itemize}
  \item 单 token：短 = 安全但每天登录 / 长 = 体验好但被偷一辈子能用。\textbf{双 token 把安全与体验两个 KPI 解耦}：
  \item \texttt{access} 1 天 (安全) + \texttt{refresh} 10 天 (留存)；每次请求都刷新一次 token 对——用户 10 天内来一次永远不过期。
  \item 10 天就是"我们能接受多久召回不到"的业务参数。
\end{itemize}
\srcnote{consts/user.go:31--32 1d / 10d；middleware/jwt.go:32 ParseRefreshToken；:59 SetToken；pkg/utils/jwt/jwt.go:81--89 续期}
\end{frame}

\begin{frame}{10 天这条线：留存与安全的平衡}
\begin{itemize}
  \item \textbf{业务侧问题}：用户多久不开 app 算流失？电商场景 7--14 天的"周期性消费"是核心，gomall 选 10 天覆盖一整个"双周购物习惯"。
  \item 30 天 = 留存好但被盗号窗口变长；3 天 = 安全但每周登录一次——\textbf{产品决定，不是开发决定}。
\end{itemize}
\begin{tabular}{@{}lll@{}}
\toprule
被盗场景 & 攻击者能用多久 & 业务损失 \\
\midrule
access 被嗅探   & 最长 1 天          & 限制在窗口内 \\
access + refresh 都泄露 & 最长 10 天 & 整个续期周期 \\
设备物理被偷    & 改密 / 主动下线之前 & 看发现速度 \\
\bottomrule
\end{tabular}
\vspace{1mm}
\small 真问题：gomall 现在没有 token 黑名单——改密码后旧 access 仍 1 天内有效。改进见"遗留风险"。
\srcnote{consts/user.go:32 RefreshTokenExpireDuration = 10 * 24 * time.Hour；middleware/jwt.go:32--63 未查黑名单}
\end{frame}

% ============================================================
\section{admin/bootstrap：开机钥匙}
% ============================================================

\begin{frame}{冷启动悖论：第一个 admin 从哪里来？}
\begin{itemize}
  \item \texttt{users/promote} 把人升级为 admin。但这个接口本身要求调用者是 admin。
  \item 第一天的生产数据库一个 admin 都没有。
  \item 这是一个鸡生蛋问题：
  \begin{itemize}
    \item 不能让普通用户能 promote 自己 $\to$ 全网随便提权。
    \item 不能直接改数据库 $\to$ DBA 权限滥用 / 没有审计。
    \item 不能写死 super-admin 账号 $\to$ 密码外泄一次完蛋。
  \end{itemize}
  \item 答案：留一个"一次性后门"，并且后门自己知道"用过一次就锁"。
\end{itemize}
\srcnote{service/admin.go:66--80 BootstrapPromoteSelf}
\end{frame}

\begin{frame}{bootstrap 状态机：用一次就锁}
\begin{tikzpicture}[node distance=10mm and 16mm]
  \node[node-state,node-state-normal] (s0) {初始\\无 admin};
  \node[node-state,node-state-transition,right=of s0] (s1) {首次调用\\count(admin)=0};
  \node[node-state,node-state-alert,right=of s1] (s2) {已锁\\count(admin) $\ge$ 1};
  \draw[arrow]       (s0) -- node[above,font=\tiny]{POST /admin/bootstrap} (s1);
  \draw[arrow]       (s1) -- node[above,font=\tiny]{Role = admin} (s2);
  \draw[arrow-bad,bend right=30] (s2.south) to node[below,font=\tiny]{再次调用 $\to$ 直接拒绝} (s0.south);
\end{tikzpicture}
\vspace{3mm}
\begin{itemize}
  \item 仍然要求"已登录"，所以不是完全裸奔。
  \item 校验 \texttt{count(role='admin') == 0}，否则报错"系统已存在 admin，禁止使用 bootstrap"。
  \item 第一个登录用户调一次 $\to$ 自己变成 admin $\to$ 之后所有人调都被拒。
\end{itemize}
\srcnote{service/admin.go:72--78 count(admin) 校验；routes/router.go:122 仅 authed，未挂 RequireRole}
\end{frame}

\begin{frame}{业务约束 (写给运维 / 安全)}
\begin{itemize}
  \item \textbf{生产 SOP}：上线 $\to$ 注册第一个运营账号 $\to$ 调 \texttt{admin/bootstrap} $\to$ user\_id 记入交接文档。之后提权走 \texttt{users/promote}（有 admin 操作日志）。
  \item \textbf{DB 清空重建}：bootstrap 又解锁；备份恢复 SOP 必须包含"恢复后立刻确认有 admin"。
  \item \textbf{不加 IP 白名单 / 时间窗口}：那是"安全感"不是"安全"，反而让人忘记基本约束。
  \item 想更硬可以要求环境变量带一次性 token，但当前业务量没必要。
\end{itemize}
\srcnote{routes/router.go:122 仅在 authed 下，无 RequireRole}
\end{frame}

% ============================================================
\section{RBAC 30 秒缓存：换 QPS 换出来的体验}
% ============================================================

\begin{frame}{为什么 30s 缓存：换 QPS 换出来的体验}
\begin{tikzpicture}[node distance=8mm and 12mm]
  \node[node-box] (req)  {admin 请求};
  \node[node-box,right=of req,node-warn] (cache) {sync.Map\\TTL 30s};
  \node[node-box,below=of cache] (db) {MySQL\\user.role};
  \node[node-box,right=of cache,node-ok] (ok) {放行 / 拒绝};
  \draw[arrow] (req) -- (cache);
  \draw[arrow,dashed] (cache) -- node[right,font=\tiny]{miss} (db);
  \draw[arrow,dashed] (db.east) to [bend right=20] node[right,font=\tiny]{回填} (cache.east);
  \draw[arrow] (cache) -- (ok);
\end{tikzpicture}
\vspace{2mm}
\begin{itemize}
  \item 鉴权每个 admin 请求都走——每次 \texttt{SELECT role FROM user} 会 QPS 翻倍 + 连接池竞争。
  \item 角色变更频率 < 一年 10 次。"30s 最终一致"换"DB 不被压垮"，业务可接受。
  \item 命中微秒 / 未命中回填 / 提权降权主动调 \texttt{InvalidateRoleCache}，客服侧"立刻生效"仍是秒级。
\end{itemize}
\srcnote{middleware/rbac.go:24,28--45,88--91；service/admin.go:59 提权后调 Invalidate}
\end{frame}

\begin{frame}{30s 不一致 + 链路开销实测}
\begin{tabular}{@{}lll@{}}
\toprule
事件 & 30s 内观感 & 是否可接受 \\
\midrule
A 被提权为 admin & A 看不到后台菜单    & 让他刷新；可接受 \\
B 被降权为普通   & B 仍能访问 admin 30s & \textbf{真风险，显式失效} \\
C 被封号         & C 仍能下单 30s       & 严重，需 token 黑名单 \\
\bottomrule
\end{tabular}
\vspace{2mm}
\small \textbf{对称性陷阱}：提权是体验问题，降权 / 封号是安全问题——gomall 用 \texttt{InvalidateRoleCache} 补这条线。但 access 1 天 TTL 仍是遗留风险。
\vspace{2mm}
\begin{tabular}{@{}lrr@{}}
\toprule
压测对照 & RPS & p95 \\
\midrule
\texttt{/ping}（无 Auth）           & 64{,}254 & 3.51ms \\
\texttt{/orders/list}（带 Auth）    & 58{,}406 & 5.00ms \\
\bottomrule
\end{tabular}
\vspace{1mm}
\small AuthMiddleware 损耗 < 10\%，\textbf{不是性能瓶颈}——继续优化先看 SQL / 分页。
\srcnote{middleware/rbac.go:88--91 显式失效；stressTest/REPORT.md \S 链路压测表 行 1, 4}
\end{frame}

% ============================================================
\section{客服补券 vs 黑产温床：边界在哪}
% ============================================================

\begin{frame}{补券这件事必须有两条路径}
\begin{itemize}
  \item \textbf{用户主路径}：\texttt{coupon/claim} $\to$ 排队 + Lua 单原子防超发 + 滑动窗口防脚本。
  \item \textbf{客服主路径}：客服核对用户身份 $\to$ 走 admin 接口直接发券到指定账户。
  \item 两条路径的本质区别：
  \begin{itemize}
    \item 用户路径是\textbf{机器对机器}的接口，没有"人核对身份"这一步 $\to$ 必须靠风控。
    \item 客服路径有\textbf{真人}核对身份 $\to$ 风控可以松，但操作要有审计。
  \end{itemize}
  \item \textbf{绝对不能做的事}：把"用户自助申请补券"做成一个普通用户能直接调的接口。
\end{itemize}
\srcnote{routes/router.go:70 coupon/claim 在 authed；admin 侧补券需新增 admin handler}
\end{frame}

\begin{frame}{为什么"用户自助补券" = 黑产温床}
\begin{tikzpicture}[node distance=8mm and 14mm]
  \node[node-box,node-state-normal] (l1)  {L1 网络层\\IP TokenBucket};
  \node[node-box,node-state-normal,right=of l1] (l2) {L2 业务层\\用户 SlidingWindow};
  \node[node-box,node-state-alert,right=of l2] (l3) {L3 客服\\真人核对};
  \draw[arrow] (l1) -- (l2);
  \draw[arrow-alert] (l2) -- (l3);
\end{tikzpicture}
\vspace{3mm}
\begin{itemize}
  \item L1 + L2 能挡 80\% 黑产，但是\textbf{挡不住有耐心的羊毛党}：注册 1 万个号，每号每天领 5 张。
  \item "用户自助补券"如果做成接口 $\to$ \textbf{第三层风控直接消失}。L1 + L2 那点防御不够。
  \item 借鉴 deck 10 的分层结论：电商场景的最后一道闸 \textbf{必须是人}，不能是代码。
\end{itemize}
\vspace{1mm}
\small \textbf{客服走 admin 的约束}：身份核对在前（工单 / 录音）；操作走 admin 接口直接发指定 user\_id；留痕（谁/何时/给谁/凭据）；单客服日额度封顶 + 异常激增告警。这套在 gomall 还没完整落地，是\textbf{下一步要做的事}。
\srcnote{deck 10 L1--L4 拦截层级；routes/router.go:22 IP 桶 + :112--119 用户滑动窗口；:125--131 admin 子组当前只有 promote / backfill}
\end{frame}


% ============================================================
\section{业务码与客服话术}
% ============================================================

\begin{frame}{鉴权失败的业务码决策树}
\begin{tikzpicture}[node distance=8mm and 10mm]
  \node[node-box] (req) {请求到达};
  \node[node-box,right=of req,node-warn] (tok) {有 token?};
  \node[node-box,right=of tok,node-warn] (parse) {签名合法?};
  \node[node-box,right=of parse,node-warn] (role) {角色匹配?};
  \node[node-box,below=of parse,node-bad] (e1) {30001\\Token 鉴权失败};
  \node[node-box,below=of role,node-bad]  (e2) {30005\\权限不足};
  \node[node-box,below=of tok,node-bad]   (e0) {400\\token 不能为空};
  \node[node-box,right=of role,node-ok]   (ok) {200\\放行};
  \draw[arrow] (req) -- (tok);
  \draw[arrow] (tok) -- (parse);
  \draw[arrow] (parse) -- (role);
  \draw[arrow] (role) -- (ok);
  \draw[arrow-bad] (tok) -- (e0);
  \draw[arrow-bad] (parse) -- (e1);
  \draw[arrow-bad] (role) -- (e2);
\end{tikzpicture}
\vspace{2mm}
\srcnote{middleware/jwt.go:22--57 三类失败分支；middleware/rbac.go:75--82 角色不匹配}
\end{frame}

\begin{frame}{客服话术：业务码 $\to$ 提示文案}
\begin{tabular}{@{}lll@{}}
\toprule
状态码 & 含义       & 客服话术 \\
\midrule
30001  & Token 鉴权失败  & "请退出重新登录"          \\
30002  & Token 已超时    & "登录已过期，请重新登录"   \\
30005  & 权限不足        & "该功能仅限管理员"        \\
60001  & 幂等 token 无效 & "请重新进入下单页"         \\
60002  & 幂等命中重复    & "您已下单成功，无需重复"   \\
70001  & 限流命中        & "活动太火爆，请稍后重试"   \\
70002  & 熔断打开        & "支付服务暂忙，1 分钟后"   \\
\bottomrule
\end{tabular}
\vspace{2mm}
\small 30001 是 token 被改 / 没带；30002 是过期；30005 要警惕用户走错入口或脚本试 admin。\\
所有鉴权失败统一返 \texttt{HTTP 200 + \{status:3000x\}}（老 app 网关吞不下 401/403），\textbf{监控告警基于 status，不是 HTTP code}。
\srcnote{pkg/e/code.go:35--58；middleware/jwt.go:24,38,51 三处 c.JSON(200,...)}
\end{frame}


% ============================================================
\section{各角色看鉴权这件事}
% ============================================================

\begin{frame}{遗留风险与改进方向}
\textbf{目前的洞}：
\begin{itemize}
  \item \textbf{无 token 黑名单}：改密 / 主动下线后旧 access\_token 仍 1 天内有效。
  \item \textbf{无 admin 操作审计表}：promote / backfill 只写 log，未入审计 DB。
  \item \textbf{admin 没有二次确认}：promote 一个用户不需要主管复核。
  \item \textbf{admin 补券接口未实现}：当前客服只能通过改 DB 这种"野路子"。
\end{itemize}
\vspace{2mm}
\textbf{改进顺序}（按 ROI）：admin 补券接口（业务最痛）$\to$ admin 操作审计表（合规）$\to$ token 黑名单（被盗号止血）$\to$ admin 二次确认（额度更大才需要）。
\srcnote{middleware/jwt.go 当前无黑名单查询；routes/router.go:125--131 admin 子组无审计中间件}
\end{frame}

\begin{frame}{回到业务：边界由谁定 + 最终目标}
\begin{tabular}{@{}ll@{}}
\toprule
决定 & 由谁负责 \\
\midrule
哪个接口能匿名访问           & 产品 + 安全 \\
access / refresh 过期时长   & 产品（留存）+ 安全（盗号窗口）\\
admin 能做哪些事             & 运营 + 安全（最小权限）\\
RBAC 缓存 TTL                & SRE + 安全（一致性 vs QPS）\\
客服补券额度上限             & 风控 + 财务 \\
bootstrap 用过谁来确认       & 运维 SOP \\
\bottomrule
\end{tabular}
\vspace{1mm}
\small \textbf{业务侧最终目标}：用户感受不到登录 / 客服 1 分钟补券 / 黑产领不到券 / 内鬼搞不出大损失（审计待补）/ 出事能止血（token 黑名单待补）。
\keypoint{鉴权不是把门关上，是\textbf{给不同角色留刚刚好够用的钥匙}。}
\end{frame}


% ============================================================
\appendix
% ============================================================

\begin{frame}{Q\&A（上）：产品 / 客服角度}
\textbf{Q1}：refresh\_token 为什么是 10 天，不是 30 天？\\
\hspace{4mm}\small A：电商双周购物周期 + 盗号窗口控制。30 天体验更好但盗号能用一个月。

\vspace{2mm}
\textbf{Q2}：客服怎么强制下线一个被盗号？\\
\hspace{4mm}\small A：现在做不到。需要补 token 黑名单。临时方案：改密 + 让用户主动等 1 天。这是当前最大的遗留风险。

\vspace{2mm}
\textbf{Q3}：用户改了密码，旧 access\_token 怎么办？\\
\hspace{4mm}\small A：当前仍然 1 天内有效。改进方向：改密 $\to$ 写黑名单 $\to$ AuthMiddleware 查黑名单。

\vspace{2mm}
\textbf{Q4}：为什么不允许"用户自助申请补券"？\\
\hspace{4mm}\small A：补券路径如果没有"人核对身份"这一步，等价于把券免费送给黑产。
\end{frame}

\begin{frame}{Q\&A（下）：安全 / SRE 角度}
\textbf{Q5}：为什么 RBAC 用 30s 缓存，不是强一致每次查 DB？\\
\hspace{4mm}\small A：高 QPS 下 DB 扛不住，30s 缓存 + 显式失效是性价比最高的方案。降权 / 封号显式调 \texttt{InvalidateRoleCache} 保证关键路径快速生效。

\vspace{2mm}
\textbf{Q6}：admin/bootstrap 会不会被攻击？\\
\hspace{4mm}\small A：要求"已登录"+"系统无 admin"两个条件同时满足才能调用。生产首发后立刻有 admin，接口自然锁死。如果数据库被清空，应当先恢复 admin 再开 API。

\vspace{2mm}
\textbf{Q7}：鉴权链路压测下会不会成为瓶颈？\\
\hspace{4mm}\small A：\texttt{orders/list} 带 Auth 仍跑到 58K RPS，相对裸 ping 损耗 < 10\%。当前不是瓶颈。

\vspace{2mm}
\textbf{Q8}：30001 大量出现意味着什么？\\
\hspace{4mm}\small A：要么 JWT secret 刚换过，要么有人在批量试 token（撞库）。后者要立刻看 IP 分布。
\end{frame}

\begin{frame}{代码位置一览}
\begin{description}
  \item[authed group 入口]\srcnote{routes/router.go:45--46}
  \item[admin sub-group]\srcnote{routes/router.go:125--131}
  \item[admin/bootstrap 路由]\srcnote{routes/router.go:122}
  \item[AuthMiddleware]\srcnote{middleware/jwt.go:16--64}
  \item[双 token 续期]\srcnote{pkg/utils/jwt/jwt.go:67--93}
  \item[RequireRole + 30s 缓存]\srcnote{middleware/rbac.go:24, 28--45, 48--86}
  \item[InvalidateRoleCache]\srcnote{middleware/rbac.go:88--91}
  \item[BootstrapPromoteSelf]\srcnote{service/admin.go:64--80}
  \item[Role 字段定义]\srcnote{repository/db/model/user.go:23, 30--31}
  \item[业务码定义]\srcnote{pkg/e/code.go:35--58；pkg/e/msg.go:31--52}
\end{description}
\small 配套压测数据：stressTest/REPORT.md 链路压测表前 7 行。
\end{frame}

\end{document}
